\begin{definition}[Canonical local-constant realization]
\label{mod:delta-localconstantrealization}
Let $F$ be a nonarchimedean local field.  Define a total canonical function
\[
 \texttt{localConstantFunction}(F):
 \texttt{LocalConstantFunction}(F)
 =\bigl(\chi:F^\times\to\mathbf C^\times\bigr)
  \longmapsto
 \bigl(\psi:F\to\mathbf C^\times\bigr)
  \longmapsto\mathbf C.
\]
Retain \texttt{localConstant}(F) as the definitionally equal short exported
name used by the final First Main Lemma.
If $\psi\ne1$, use Theorem~\ref{mod:basic-characterconductorexistence} to
form the canonical exact packages for $\chi$ and $\psi$, choose an admissible
denominator $\gamma$, and set the value equal to
\[
 \Delta_F^{\mathrm{fin}}(\chi,\psi;\gamma)
 =\texttt{deltaFinite}(\chi,\psi,\gamma).
\]
Independence of $\gamma$ is the existing theorem from
Definition~\ref{mod:delta-finitedefinition}; the definition must not retain a
mathematically observable choice of conductor proof or denominator.

When $\psi=1$, set the value to zero only to make
\texttt{localConstantFunction} total.  This branch is outside the manuscript's
nontrivial-additive-character local constant and must not be used to assert a
finite or integral formula that requires exact additive-conductor data.

Prove the computational specification
\[
 \texttt{IsDeltaFiniteLocalConstant}
   (\texttt{localConstantFunction}(F)).
\]
Then prove the finite-value realization for every ordinary continuous
quasi-character and every nontrivial ordinary continuous additive character,
with any admissible denominator for their canonical conductor packages.
Finally, for every positive left Haar measure $du$ on $U_F^0$, prove the
integral realization
\[
 \texttt{localConstantFunction}(F)(\chi)(\psi)
 =\Delta_F^{\mathrm{int}}(\chi,\psi;\gamma)
\]
by applying the existing integral--finite bridge with the same packages and
denominator.  No integral comparison is reproved in this node.
\LeanFile{LanglandsFirstMainLemma/Delta/LocalConstantRealization.lean}
\lean{LanglandsFirstMainLemma.localConstantFunction}
\lean{LanglandsFirstMainLemma.localConstantFunction_isDeltaFinite}
\lean{LanglandsFirstMainLemma.localConstantFunction_eq_deltaFinite}
\lean{LanglandsFirstMainLemma.localConstantFunction_eq_deltaIntegral}
\lean{LanglandsFirstMainLemma.localConstant}
\leanok
\SourceText{Definition of Delta; Lemma lem:gamma-independent; Corollary cor:local-gauss-integral-nonzero; final export contract.}
\uses{mod:basic-characterconductorexistence,mod:delta-errorterm,mod:delta-integralfinitebridge}
\FormalizationContract
This node is expected to occupy roughly 140--300 Lean lines once implemented.
Its finite branch must use \texttt{deltaFinite}, its computational theorem
must have the existing \texttt{IsDeltaFiniteLocalConstant} predicate, and its
integral theorem is obtained from
\texttt{deltaIntegral\_eq\_deltaFinite}.  The trivial additive-character
case is only a totalization device.
\end{definition}
